% =====================================================================
% BUSINESS CASE — MBA USP/Esalq em Engenharia de Software
% Gisele Fonseca de Aguiar Silva
%
% COMO USAR NO OVERLEAF
%   Arquivo UNICO: nao precisa de nenhum outro arquivo para compilar.
%   1. Novo projeto -> Blank Project -> apague o conteudo e cole este
%      arquivo inteiro. Ou Upload Project enviando so este .tex.
%   2. Compilador: pdfLaTeX (ja e o padrao do Overleaf).
%   3. A logo oficial do programa e OPCIONAL para compilar: baixe-a do
%      Sistema e envie como logo-mba.png. Sem ela o documento compila
%      normalmente e mostra uma caixa no lugar da imagem.
%   4. Para usar a Arial de verdade, veja a observacao no fim do bloco
%      de normas e compile com LuaLaTeX.
%
% ESTRUTURA OBRIGATÓRIA (item 6 do manual), nesta ordem:
%   Folha de Rosto -> Título + Sumário Executivo -> Introdução ->
%   Desenvolvimento -> Discussão dos Resultados -> Considerações Finais
%   -> Referências -> Apêndice(s)/Anexo(s) (opcionais)
%
% LIMITE DE PÁGINAS: 30, incluindo apêndices e anexos.
%   O manual abre exceção de 50 páginas para Business Cases de
%   Engenharia de Software com temática de Implementação de
%   Algoritmo(s) de Machine Learning. Este trabalho NÃO implementa
%   algoritmo, então o limite aplicável é 30 até que o Especialista
%   de Business Case confirme o contrário.
% =====================================================================
\documentclass[11pt,a4paper]{article}
% ================================================================
% NORMAS DE FORMATACAO MBA USP/Esalq, embutidas neste arquivo
% (Manual de Instrucoes e Normas: Tabelas 3, 4 e 5; itens 17 e 18)
% Arquivo unico de proposito: NAO depende de nenhum .sty externo.
% ================================================================
% =====================================================================
% Normas de formatação do Business Case MBA USP/Esalq
% Codificado a partir do Manual de Instruções e Normas (Tabelas 3, 4 e 5
% e itens 17 e 18). Cada regra está comentada com a sua origem.
% =====================================================================

% ---------- Codificação e idioma ----------
\usepackage[T1]{fontenc}
\usepackage[utf8]{inputenc}
% No Overleaf o idioma brasileiro está sempre disponível. O teste abaixo
% só existe para que o arquivo também compile em instalações TeX enxutas,
% sem o pacote de hifenização portuguesa.
\IfFileExists{portuges.ldf}
  {\usepackage[brazil]{babel}}
  {\usepackage[english]{babel}
   \GenericWarning{}{Hifenizacao portuguesa ausente;
     compilando em ingles. No Overleaf isto nao ocorre}}

% ---------- Fonte: Arial 11, cor preta (Tabela 3) --------------------
% Arial não acompanha o TeX. TeX Gyre Heros é clone metricamente
% compatível com Helvetica/Arial e é o substituto aceito em pdfLaTeX.
% Para usar a Arial de verdade, compile com LuaLaTeX e troque este
% bloco pelo bloco comentado ao final deste arquivo.
% Sem escala: Arial e Helvetica compartilham as mesmas métricas, então
% 11 pt aqui corresponde a Arial 11 pt. Não use [scaled], que reduziria
% o corpo e descumpriria a norma.
\usepackage[scaled=1.0]{helvet}
\renewcommand{\familydefault}{\sfdefault}

% ---------- Página: A4 retrato, margens 2,5 cm (Tabela 3) ------------
\usepackage[a4paper,top=2.5cm,bottom=2.5cm,left=2.5cm,right=2.5cm,
                headheight=38pt,headsep=12pt,footskip=1.0cm]{geometry}

% ---------- Espaçamento (Tabela 4) ----------------------------------
% 1,5 é o padrão do corpo do texto; 1,0 em seções específicas.
\usepackage{setspace}
\onehalfspacing

% ---------- Recuo de primeira linha: 1,25 cm (Tabela 3) -------------
\setlength{\parindent}{1.25cm}
\setlength{\parskip}{0pt}

% ---------- Texto justificado (Tabela 3) ----------------------------
% É o padrão do LaTeX. Referências ficam alinhadas à esquerda — ver
% o ambiente referencias mais abaixo.
\usepackage{ragged2e}

% ---------- Cabeçalho repetido e paginação (item 18.1) --------------
% Cabeçalho em todas as páginas, alinhado à esquerda, com a logo no
% canto superior direito. Paginação em arábicos no canto inferior
% direito, Arial 9, iniciando na folha de rosto (página 1).
\usepackage{fancyhdr}
\usepackage{graphicx}

% Sem nomes com @, para que este bloco também funcione quando colado
% direto no preâmbulo de um .tex, onde @ não é letra.
\newcommand{\bccursonome}{Engenharia de Software}
\newcommand{\bcanodefesa}{2027}
\newcommand{\bccurso}[1]{\renewcommand{\bccursonome}{#1}}
\newcommand{\bcano}[1]{\renewcommand{\bcanodefesa}{#1}}

\newcommand{\bccabecalho}{%
  \fontsize{9}{11}\selectfont
  \begin{minipage}[c]{0.74\textwidth}
    \RaggedRight Business Case apresentado para obtenção do título de
    especialista em \bccursonome{} -- \bcanodefesa{}
  \end{minipage}%
  \hfill
  \begin{minipage}[c]{0.24\textwidth}
    \raggedleft
    % A logo oficial do programa deve ser baixada do Sistema e salva
    % como logo-mba.png na raiz do projeto. Sem o arquivo, o documento
    % compila mesmo assim e mostra um aviso no lugar.
    \IfFileExists{logo-mba.png}{%
      \includegraphics[height=1.05cm]{logo-mba.png}}{%
      \fbox{\scriptsize\ttfamily logo-mba.png}}
  \end{minipage}%
}

\fancypagestyle{bcpadrao}{%
  \fancyhf{}
  \renewcommand{\headrulewidth}{0pt}
  \renewcommand{\footrulewidth}{0pt}
  \fancyhead[L]{\bccabecalho}
  \fancyfoot[R]{\fontsize{9}{11}\selectfont\thepage}
}
\pagestyle{bcpadrao}

% ---------- Títulos de seção (Tabela 3 e itens 18.3 a 18.7) ---------
% Negrito, alinhados à esquerda, sem recuo e SEM numeração.
% Separados do texto por um espaço de caractere.
\usepackage[explicit]{titlesec}

\titleformat{\section}{\normalfont\bfseries}{}{0pt}{\RaggedRight #1}
\titlespacing*{\section}{0pt}{\baselineskip}{\baselineskip}

% Subtítulos: negrito, primeiras letras maiúsculas, sem numeração,
% com recuo especial de 1,25 cm. Usar apenas quando fundamentais.
\titleformat{\subsection}[hang]{\normalfont\bfseries}{}{0pt}
            {\hspace*{1.25cm}#1}
\titlespacing*{\subsection}{0pt}{\baselineskip}{\baselineskip}

% Impede numeração automática em qualquer nível.
\setcounter{secnumdepth}{0}

% ---------- Figuras e Tabelas (itens 17.1 e 17.2) -------------------
% Títulos, fontes e notas em espaçamento simples. O uso de QUADROS
% não é permitido pelo manual — só Figuras e Tabelas.
\usepackage{caption}
\usepackage{booktabs}
\usepackage{array}

% Corpo 11 como no texto (a norma não prevê legenda em corpo menor).
% Os nomes são fixados à mão para não dependerem do idioma do babel.
\captionsetup{font={sf},labelsep=period,
              justification=justified,singlelinecheck=false,
              skip=4pt}
\captionsetup[table]{position=above,name=Tabela}
\captionsetup[figure]{position=below,name=Figura}
% Confira o estilo do rótulo contra o exemplo dos Anexos do manual: as
% tabelas do próprio manual aparecem como "Tabela 3." sem negrito. Para
% usar negrito, acrescente labelfont=bf acima.

% Fonte obrigatória logo abaixo da figura/tabela, espaçamento simples.
\newcommand{\fonteprop}[1]{%
  \par\vspace{2pt}%
  {\singlespacing\fontsize{11}{13}\selectfont\RaggedRight Fonte: #1\par}}
\newcommand{\fonteoriginal}{\fonteprop{Resultados originais do estudo}}

% Nota de tabela, quando houver, vem APÓS a fonte.
\newcommand{\notatab}[1]{%
  \par{\singlespacing\fontsize{11}{13}\selectfont\RaggedRight Nota: #1\par}}

% ---------- Sumário Executivo (item 18.2) ---------------------------
% Parágrafo único, SEM recuo na primeira linha, espaçamento simples,
% no máximo 250 palavras, pretérito perfeito, sem citações.
\newenvironment{sumarioexecutivo}
  {\par\vspace{\baselineskip}%
   \noindent\textbf{Sumário Executivo}\par\vspace{\baselineskip}%
   \begingroup\singlespacing\setlength{\parindent}{0pt}}
  {\par\endgroup\vspace{\baselineskip}}

% ---------- Referências (item 18.7) ---------------------------------
% Alinhadas à esquerda, sem recuo de primeira linha, espaçamento
% simples. Nenhum estilo padrão do LaTeX reproduz os modelos do item 21
% do manual, então a lista é escrita à mão, o que também garante
% aderência. Use \refitem para cada entrada.
\newenvironment{referencias}
  {\begingroup\singlespacing\RaggedRight\setlength{\parindent}{0pt}}
  {\par\endgroup}
\newcommand{\refitem}[1]{#1\par\vspace{0.5\baselineskip}}

% ---------- Normatização do texto (Tabela 5) ------------------------
% Sigla entre COLCHETES após a definição, na primeira ocorrência.
\newcommand{\sigla}[2]{#1 [#2]}
% Expressão estrangeira não corrente vai entre aspas; latim em itálico.
\newcommand{\estrangeiro}[1]{``#1''}
\newcommand{\latim}[1]{\emph{#1}}
% Separador de milhar com ponto e decimal com vírgula.
\usepackage{siunitx}
\sisetup{output-decimal-marker={,},group-separator={.},
         group-minimum-digits=4,detect-all}

% ---------- Utilidades ----------------------------------------------
\usepackage{xcolor}
\usepackage[hidelinks]{hyperref}
\usepackage{microtype}

% Marcador visível de pendência. Rode \pendenciasoff antes de depositar
% para conferir que não sobrou nenhuma.
\newcommand{\pend}[1]{\textcolor{red}{\textbf{[PENDENTE: #1]}}}
\newcommand{\pendenciasoff}{\renewcommand{\pend}[1]{}}

% =====================================================================
% Para Arial de verdade: troque o bloco de fonte acima por este e
% compile com LuaLaTeX (menu Overleaf: Compilador -> LuaLaTeX).
% A Arial precisa ser enviada como arquivo .ttf para o projeto.
%
% \usepackage{fontspec}
% \setmainfont{Arial}[
%   BoldFont       = Arial Bold,
%   ItalicFont     = Arial Italic,
%   BoldItalicFont = Arial Bold Italic]
% =====================================================================
% ----------------- fim das normas -------------------------------

% Curso e ano de defesa, usados no cabeçalho de todas as páginas.
\bccurso{Engenharia de Software}
\bcano{2027}

\begin{document}

% =====================================================================
% FOLHA DE ROSTO — página 1 (item 18.1)
% Espaçamento simples. Título em negrito, centralizado, no máximo 15
% palavras, com maiúscula apenas na inicial e em nomes próprios.
% =====================================================================
\begin{singlespace}
\thispagestyle{bcpadrao}
\vspace*{1.5cm}

\begin{center}
\textbf{Custo do erro silencioso em pipelines analíticos: onde posicionar a
validação de dados}
\end{center}

\vspace{2\baselineskip}

\begin{center}
Gisele Fonseca de Aguiar Silva\textsuperscript{1}*;
\pend{Nome do Especialista de Business Case}\textsuperscript{2}
\end{center}

\vspace{\baselineskip}

{\fontsize{9}{11}\selectfont\RaggedRight
\noindent\textsuperscript{1}Mestra em Administração. Engenheira de dados.
E-mail do autor correspondente: gih52sim@gmail.com\par
\noindent\textsuperscript{2}\pend{Titulação do Especialista}.
\pend{Instituição}. E-mail: \pend{e-mail}\par
}
\end{singlespace}

\newpage

% =====================================================================
% TÍTULO + SUMÁRIO EXECUTIVO — página 2 (item 18.2)
% Arial 11, espaçamento simples. Parágrafo único, sem recuo, no máximo
% 250 palavras, pretérito perfeito do indicativo, SEM citações.
% =====================================================================
\begin{singlespace}
\begin{center}
\textbf{Custo do erro silencioso em pipelines analíticos: onde posicionar a
validação de dados}
\end{center}
\end{singlespace}

\begin{sumarioexecutivo}
\pend{redigir ao final, quando os resultados estiverem fechados --- 250
palavras, parágrafo único, pretérito perfeito, sem citações. Rascunho
provisório abaixo, a revisar.}
Pipelines analíticos entregaram métricas incorretas sem gerar qualquer
alarme quando o código de transformação continha defeito semântico, isto
é, lógica que executou normalmente mas não correspondeu à intenção de
negócio. O estudo examinou esse problema no contexto de equipes de
engenharia de dados que mantêm transformações que não escreveram.
Construiu-se um experimento de injeção de defeitos sobre um pipeline
analítico com dados simulados gerados a partir de procedimentos
estatísticos declarados, no qual nove mecanismos de defeito, ancorados em
taxonomia publicada, foram confrontados com quatro camadas de controle de
qualidade de dados, medindo-se o erro induzido em sete métricas de
negócio ao longo de trinta repetições. \pend{complementar com a segunda
frente (levantamento de projetos públicos) e com os resultados finais.}
Os testes de schema e os testes genéricos de unicidade, nulidade e
integridade referencial não detectaram nenhum dos nove defeitos, por
verificarem a forma da tabela e não o significado do número. A
reconciliação de contagem de linhas reduziu a exposição de forma
expressiva, mas permaneceu cega aos defeitos que preservam a contagem, um
dos quais produziu erro superior a sessenta por cento no faturamento
apurado. O estudo entregou uma matriz de alocação de controles, indicando
qual verificação posicionar em cada camada conforme o mecanismo de
defeito e as métricas consumidas pelo negócio, com o risco residual de
cada configuração dimensionado em valor financeiro.
\end{sumarioexecutivo}

% =====================================================================
% INTRODUÇÃO (itens 6.2 e 18.3)
% No máximo DUAS páginas (~35 linhas por página). Proibido: subtópicos,
% marcadores, figuras, tabelas e quadros. Três partes, nesta ordem:
%   (i) problema ou oportunidade; (ii) contexto; (iii) objetivo.
% O objetivo deve ser enunciado no ÚLTIMO parágrafo.
% =====================================================================
\section{Introdução}

\pend{Parte (i) --- problema ou oportunidade. Fórmula sugerida pelo
material do curso: ``o problema de A tem o efeito de B, com o impacto de
C''. Aqui: código de transformação que não foi escrito por quem o mantém
produz erro semântico que executa em silêncio (A), a métrica entregue ao
negócio sai errada sem alarme (B), e a decisão tomada sobre ela carrega
exposição financeira não medida (C).}

\pend{Parte (ii) --- contexto. Engenharia de dados em ambiente de setor
financeiro; a prática corrente de validação apoiada em testes genéricos e
de schema; e a literatura recente sobre falha silenciosa em SQL gerado
por modelos de linguagem, que documenta a taxa de erro semântico mas
avalia a consulta isolada, e não o dano à métrica quando a consulta vira
modelo de transformação em produção.}

\pend{Parte (iii) --- objetivo, no último parágrafo, único, claro,
específico e mensurável. Dizer o que se vai MEDIR, não ``analisar'' nem
``compreender''. Rascunho: ``O objetivo deste trabalho foi quantificar o
erro induzido em métricas de negócio por mecanismos de defeito semântico
em modelos de transformação de dados e determinar a capacidade de
detecção de cada camada de controle de qualidade, estabelecendo uma
matriz de alocação de controles com o risco residual de cada configuração
expresso em valor financeiro.''}

% =====================================================================
% DESENVOLVIMENTO (itens 6.3 e 18.4)
% Subtítulos são permitidos e ajudam. Figuras e Tabelas permitidas,
% QUADROS não. Conteúdo esperado: caracterização do objeto de estudo,
% diagnóstico do problema, métodos aplicados e, quando aplicável, plano
% de implementação. Em trabalho com viés científico, esta seção descreve
% a metodologia: tipo de pesquisa, abordagem, fontes, instrumentos de
% coleta, critérios de análise e procedimentos de validade.
% =====================================================================
\section{Desenvolvimento}

\subsection{Caracterização do Objeto de Estudo}

O objeto deste estudo é a camada de transformação de pipelines analíticos
e o conjunto de controles de qualidade que a cerca. Entende-se por camada
de transformação o código que converte dados brutos ingeridos de sistemas
de origem em tabelas agregadas, consumidas por relatórios e indicadores de
negócio. Em arquiteturas correntes essa camada é escrita predominantemente
em \estrangeiro{Structured Query Language} [SQL] e organizada em modelos
versionados, executados em cadeia por um orquestrador.

O ambiente de referência é o de equipes de engenharia de dados que operam
pipelines em produção no setor financeiro, no qual indicadores derivados
alimentam decisões de crédito, exposição e prestação de contas
regulatória. Três características desse ambiente delimitam o problema
tratado. A primeira é a separação entre quem escreveu a transformação e
quem a mantém: parte relevante da lógica em operação foi construída por
equipes anteriores e é sustentada por profissionais que não participaram
da sua concepção. A segunda é a tolerância a falhas adotada por projeto,
em que conversões e junções são deliberadamente permissivas para que a
execução não interrompa em presença de dado inesperado. A terceira é a
assimetria entre a velocidade de produção de código de transformação e a
velocidade de produção de verificação, acentuada pela adoção de
assistentes de geração de código.

As partes interessadas são a equipe de engenharia de dados, responsável
pela construção e manutenção dos modelos; as áreas de negócio, que
consomem os indicadores sem acesso à lógica que os produz; e as funções de
controle interno e auditoria, que respondem pela fidedignidade dos números
reportados. A assimetria de informação entre esses grupos é o que torna a
falha silenciosa relevante: o consumidor do indicador não dispõe de meio
próprio para questioná-lo.

\subsection{Diagnóstico do Problema}

O defeito tratado neste estudo não é a interrupção do pipeline, e sim a
sua conclusão bem-sucedida com resultado incorreto. Um modelo de
transformação com defeito semântico é aquele cuja consulta é sintaticamente
válida, executa sem exceção e materializa a tabela de destino, mas não
corresponde à intenção de negócio que deveria implementar. A literatura da
área de geração automática de consultas denomina essa classe de ocorrência
falha silenciosa, por oposição à falha ruidosa, que interrompe a execução e
é, por isso, barata de detectar.

A magnitude do fenômeno está documentada para consultas geradas por
modelos de linguagem. Erros de natureza semântica respondem por parcela
expressiva do total de erros observados em conjuntos de avaliação
consagrados, e a maior parte deles decorre de má interpretação do esquema
do banco ou da pergunta formulada, e não de erro de sintaxe
\pend{inserir citações autor-ano após confirmar as autorias completas nas
fontes originais}. Trabalhos recentes de verificação relatam, ainda, que a
autoavaliação do próprio modelo gerador tende à confiança excessiva, com
baixa revocação na identificação dos seus erros, o que inviabiliza
delegar a verificação ao mesmo mecanismo que produziu o código.

Dois limites dessa literatura delimitam a lacuna aqui tratada. O primeiro
é de objeto: a avaliação recai sobre a consulta isolada, em tarefa de
pergunta e resposta, e não sobre o modelo de transformação em operação,
cujo resultado incorreto é materializado em tabela e consumido por um
indicador. O segundo é de desfecho: mede-se a correção da consulta, e não
o dano causado à métrica de negócio. Em paralelo, a literatura de
qualidade de pipelines de dados oferece taxonomia empiricamente
fundamentada de fatores de influência e causas-raiz, mas não quantifica o
custo de detectar a falha em cada camada, conforme declarado pelos
próprios autores (Foidl et al., 2024).

A restrição prática que sustenta a relevância da intervenção é a natureza
dos controles usualmente adotados. Os testes de uso corrente verificam
propriedades da forma da tabela, como unicidade da chave, ausência de
nulos, integridade referencial e conformidade de tipos. Defeito semântico,
por definição, preserva a forma: a tabela resultante continua tendo chave
única, colunas não nulas e referências válidas, ainda que os valores
estejam errados. Decorre daí a hipótese que o estudo submete a teste: a
pilha de controles de maior penetração na prática é estruturalmente
incapaz de detectar a classe de defeito que mais cresce em frequência.

\subsection{Procedimentos Metodológicos}

O estudo é de natureza aplicada e abordagem quantitativa, e combina dois
delineamentos complementares. O primeiro é um experimento controlado de
injeção de defeitos, conduzido sobre um pipeline analítico com verdade
conhecida por construção, no qual se manipula o mecanismo de defeito e se
observa, simultaneamente, o erro induzido em cada métrica de negócio e a
detecção obtida por cada camada de controle. O segundo é um levantamento
de caráter exploratório sobre projetos públicos de transformação de dados,
destinado a estabelecer quais camadas de controle são efetivamente
declaradas na prática.

A articulação entre os dois delineamentos é o que permite responder à
pergunta de gestão. O experimento estabelece o que cada camada de controle
detecta; o levantamento estabelece quais camadas as equipes de fato
adotam. O cruzamento dos dois resultados fornece a exposição associada à
configuração de controles mais frequente, que é o insumo da decisão de
alocação.

A opção por dados simulados no experimento decorre de exigência do próprio
desenho: a mensuração do erro induzido pressupõe conhecer o valor correto
da métrica, o que inexiste em base operacional. As normas institucionais
admitem a geração de dados simulados para os Business Cases de Engenharia
de Software, condicionada à fundamentação em procedimentos estatísticos
bem definidos e à sua descrição detalhada, requisito atendido na subseção
seguinte. Nenhum dado de organização privada foi utilizado em qualquer das
duas frentes.

\subsection{Geração dos Dados Simulados}

O cenário simulado reproduz um domínio de pedidos e itens de pedido, com
100 contas de cliente, 30.000 pedidos e janela de observação de 90 dias.
A métrica central é o faturamento líquido do período, obtida pela soma,
sobre os pedidos não cancelados, do produto entre quantidade, preço
unitário e complemento do desconto. Os procedimentos estatísticos de
geração estão na Tabela~\ref{tab:geracao}.

\begin{table}[htbp]
\centering
\caption{Procedimentos estatísticos de geração do cenário simulado}
\label{tab:geracao}
\begin{singlespace}
\begin{tabular}{@{}ll@{}}
\toprule
\textbf{Elemento} & \textbf{Procedimento} \\
\midrule
Porte da conta        & Lognormal($\mu=0$; $\sigma=0,8$), normalizado pela média \\
Alocação de pedidos   & Multinomial com probabilidade proporcional ao porte \\
Itens por pedido      & $1 + $ Poisson($\lambda=2$) \\
Quantidade por item   & $1 + $ Poisson($\lambda=1$) \\
Preço unitário        & Lognormal($\mu=3,2$; $\sigma=0,7$) escalado pelo porte \\
Desconto              & Uniforme discreta sobre $\{0; 0,05; 0,10; 0,15\}$ \\
Status do pedido      & Multinomial (cancelado 0,08; pendente 0,07; concluído 0,85) \\
Status ausente        & Bernoulli, $p = 0,030\,(1 + 0,90z)$, $z$ truncado em $[-1;3]$ \\
Duplicação de item    & Bernoulli($p=0,02$) \\
Pedido sem item       & Bernoulli($p=0,03$) \\
\bottomrule
\end{tabular}
\end{singlespace}
\fonteoriginal
\notatab{$z$ é o valor padronizado do pedido. A probabilidade crescente de
status ausente com o valor do pedido reproduz falha correlacionada ao
valor, em que o defeito atinge preferencialmente os registros de maior
magnitude.}
\end{table}

Quatro complicações foram deliberadamente embutidas nos dados de origem,
de modo que o cenário apresentasse as ambiguidades que um modelo de
transformação precisa resolver corretamente. A primeira é a ausência de
status em parte dos pedidos, decorrente de lacuna de ingestão, com
probabilidade crescente conforme o valor do pedido; trata-se de falha
correlacionada ao valor, que atinge preferencialmente os registros de
maior magnitude. A segunda é a duplicação de linhas de item por dupla
carga da origem, que exige deduplicação a jusante. A terceira é a
existência de pedidos válidos sem nenhum item, cuja receita é
legitimamente nula. A quarta é a defasagem entre a data do pedido e a data
de embarque, que cria a possibilidade de filtragem pela coluna errada.

\subsection{Validação do Instrumento}

Antes de qualquer injeção de defeito, verificou-se que o instrumento media
o que se propunha a medir. Para tanto, a verdade do cenário foi calculada
diretamente a partir do processo gerador, sem passar por consulta, e
comparada ao resultado de uma implementação de referência em SQL. As duas
apurações coincidiram nas sete métricas observadas, com diferença relativa
nula até a sexta casa decimal. Essa conferência é condição de validade
interna do experimento: sem ela, qualquer erro atribuído a um defeito
injetado poderia ser, na verdade, erro do próprio aparato de medição.

\subsection{Mecanismos de Defeito Injetados}

Foram implementados nove mecanismos de defeito, cada um correspondendo a
uma classe documentada na literatura de erros em consultas geradas e em
qualidade de pipelines, de modo que o catálogo não resultasse de escolha
arbitrária do pesquisador. Todos foram implementados como consultas
executáveis sobre as mesmas tabelas de origem, e todos produzem tabela de
destino sem levantar exceção. A Tabela~\ref{tab:defeitos} descreve os
mecanismos.

\begin{table}[htbp]
\centering
\caption{Mecanismos de defeito semântico injetados no modelo de transformação}
\label{tab:defeitos}
\begin{singlespace}
\fontsize{10}{12}\selectfont
\begin{tabular}{@{}p{0.8cm}p{6.6cm}p{7.2cm}@{}}
\toprule
\textbf{Cód.} & \textbf{Defeito injetado} & \textbf{Classe na taxonomia} \\
\midrule
D1 & Junção de pedidos a itens com soma do total do cabeçalho & Erro de grão, com multiplicação do valor pelo número de itens \\
D2 & Junção interna no lugar de junção externa à esquerda & Erro de síntese, com descarte de pedidos sem item \\
D3 & Comparação de desigualdade sobre coluna que admite nulo & Má interpretação do banco, pela lógica de três valores \\
D4 & Literal de filtro grafado em caixa divergente da origem & Valor de filtro fabricado, que torna o filtro inoperante \\
D5 & Filtragem pela data de embarque no lugar da data do pedido & Má interpretação do esquema \\
D6 & Omissão da deduplicação da dupla carga de itens & Omissão de etapa de limpeza \\
D7 & Desconsideração do percentual de desconto & Semântica de coluna incorreta \\
D8 & Produto entre preço máximo e quantidade somada do pedido & Erro de grão aritmético em pedido com vários itens \\
D9 & Intervalo fechado sobre carimbo de tempo & Erro de fronteira, que exclui o último dia da janela \\
\bottomrule
\end{tabular}
\end{singlespace}
\fonteoriginal
\notatab{A classificação apoia-se nas taxonomias de erro semântico em
consultas geradas e de causas-raiz em pipelines de dados
\pend{citar autor-ano após confirmar autorias}. Nenhum dos nove
mecanismos interrompe a execução.}
\end{table}

\subsection{Camadas de Controle Avaliadas}

Cada modelo defeituoso foi submetido a quatro camadas de controle,
dispostas em ordem crescente de custo de implantação, descritas na
Tabela~\ref{tab:controles}. A escolha das camadas reproduz a prática
corrente: testes de forma e de esquema são os de menor custo e maior
penetração, ao passo que a reconciliação de valor exige a existência de um
total de controle independente na origem, nem sempre disponível.

\begin{table}[htbp]
\centering
\caption{Camadas de controle de qualidade avaliadas e respectivos critérios de disparo}
\label{tab:controles}
\begin{singlespace}
\fontsize{10}{12}\selectfont
\begin{tabular}{@{}p{0.8cm}p{5.2cm}p{8.6cm}@{}}
\toprule
\textbf{Cód.} & \textbf{Camada} & \textbf{Critério de disparo} \\
\midrule
C1 & Reconciliação de contagem & Divergência superior a 0,1\% entre a contagem de pedidos no destino e a contagem apurada independentemente na origem \\
C2 & Testes genéricos de forma & Violação de unicidade da chave, de ausência de nulos na chave ou na medida, de integridade referencial ou de não negatividade \\
C3 & Esquema e contrato & Ausência de coluna esperada ou divergência de tipo na medida \\
C4 & Reconciliação de valor & Desvio superior a 5\% entre a razão do total do destino pelo total declarado na origem e a razão calibrada em período saudável \\
\bottomrule
\end{tabular}
\end{singlespace}
\fonteoriginal
\notatab{O total declarado na origem não é oráculo: inclui imposto e frete,
que o modelo analítico exclui por definição. Por isso a reconciliação de
valor opera por razão calibrada com banda de tolerância, e não por
igualdade. A banda de 5\% é parâmetro de projeto, não medida de campo; a
sensibilidade do resultado a esse valor foi avaliada e consta da Discussão
dos Resultados.}
\end{table}

\subsection{Métricas de Negócio Observadas}

Para cada execução registrou-se o erro relativo induzido em sete métricas,
escolhidas por representarem usos distintos do mesmo conjunto de dados: a
soma do faturamento, a média e a mediana do valor por pedido, o percentil
95 da distribuição de valores, a contagem de pedidos, o valor máximo
observado e a soma da conta de maior faturamento. A inclusão de métricas de
cauda e de concentração justifica-se por serem as que alimentam decisões de
exposição e limite, e por serem as menos cobertas pelos controles de uso
corrente.

Cada combinação de mecanismo de defeito foi executada sob 30 sementes
aleatórias distintas, totalizando 270 execuções, e os resultados são
reportados como média e desvio-padrão entre repetições. A repetição
responde a uma limitação de desenhos com semente única, nos quais um valor
nulo de erro pode decorrer de sorteio favorável e não de ausência de
efeito.

\subsection{Levantamento de Projetos Públicos de Transformação}

A segunda frente destina-se a estabelecer quais camadas de controle são
efetivamente declaradas em projetos reais. A população de interesse é a dos
repositórios públicos que contêm projeto de transformação de dados, e o
quadro amostral foi construído por busca de código sobre o arquivo de
configuração que caracteriza esse tipo de projeto. Como o serviço de busca
limita o número de resultados por consulta, o quadro foi montado por
estratificação em nove faixas de tamanho de repositório, procedimento que
supera esse limite sem recorrer a amostragem por conveniência.

Aplicaram-se cinco critérios de exclusão, nesta ordem, com registro da
contagem por critério: repositório derivado de outro, por não representar
decisão própria de teste; repositório arquivado, por não refletir prática
corrente; ausência de diretório de modelos; projeto com menos de cinco
modelos, por caracterizar exercício; e indício de material didático no nome
ou na descrição. A amostragem é aleatória simples sem reposição sobre o
quadro elegível, com semente fixa, de modo a permitir replicação.

A mensuração consiste na leitura dos arquivos de declaração de testes de
cada projeto e na contagem, por tipo, dos testes efetivamente declarados,
incluindo os contratos de esquema e os pacotes de teste adotados. Cada
projeto é então classificado na camada mais alta que possui, segundo o
mapeamento da Tabela~\ref{tab:mapeamento}, que é o que torna as duas
frentes comparáveis.

\begin{table}[htbp]
\centering
\caption{Correspondência entre camada de controle e teste observável em projeto real}
\label{tab:mapeamento}
\begin{singlespace}
\fontsize{10}{12}\selectfont
\begin{tabular}{@{}p{0.8cm}p{4.4cm}p{9.4cm}@{}}
\toprule
\textbf{Cód.} & \textbf{Camada} & \textbf{Declaração observável no projeto} \\
\midrule
C1 & Reconciliação de contagem & Testes de igualdade e de limite de contagem de linhas entre relações; testes de anomalia de volume \\
C2 & Testes genéricos de forma & Testes nativos de unicidade, ausência de nulos, valores aceitos e integridade referencial \\
C3 & Esquema e contrato & Contrato de modelo com imposição ativa, restrições de coluna e testes de conformidade de tipo \\
C4 & Reconciliação de valor & Pacotes de comparação entre relações; testes de expressão sobre agregados; testes singulares com indício de reconciliação \\
\bottomrule
\end{tabular}
\end{singlespace}
\fonteoriginal
\notatab{A identificação de reconciliação em teste singular é heurística,
baseada em padrões textuais, e superestima a camada C4. O valor obtido
para essa camada deve ser lido como limite superior.}
\end{table}

O instrumento de leitura foi aferido por dois procedimentos antes da
coleta. O primeiro consistiu em submetê-lo a casos construídos com gabarito
conhecido, cobrindo os formatos de declaração em uso, inclusive a
renomeação recente da chave de declaração de testes, que um leitor ingênuo
desconsideraria por completo. O segundo consistiu na conferência contra
arquivos de declaração de projetos reais, com comparação manual do
resultado. \pend{reportar a execução da coleta e o número de projetos
efetivamente analisados quando a frente for concluída.}

% =====================================================================
% DISCUSSÃO DOS RESULTADOS (itens 6.4 e 18.5)
% Seção OBRIGATÓRIA. Apresentar, discutir e interpretar os resultados
% obtidos ou projetados, de forma comparativa à literatura. Comparação
% antes/depois. Viabilidade em quatro fatores: operacional, financeira,
% estratégica e social. Citações permitidas aqui.
% =====================================================================
\section{Discussão dos Resultados}

\subsection{Detecção por Camada de Controle}

\pend{Tabela da matriz defeito $\times$ controle. Resultado central: os
testes genéricos e o teste de schema não detectaram nenhum dos nove
defeitos.}

\subsection{Risco Residual por Pilha de Controles}

\pend{Tabela das pilhas A a D com exposição somada. Discutir que adotar
testes genéricos sobre nada não reduziu exposição, que a reconciliação de
contagem teve o melhor retorno, e que o restante só cai com reconciliação
de valor.}

\subsection{Vulnerabilidade Dependente da Métrica}

\pend{Discutir o defeito invisível no faturamento e material na mediana,
e a cauda como pior caso.}

\subsection{Viabilidade da Proposta}

\pend{Quatro fatores: operacional (funciona na prática), financeira (faz
sentido econômico), estratégica (alinhada à visão da organização) e
social (impacto sobre as pessoas). Separar ganhos tangíveis de
intangíveis.}

% =====================================================================
% CONSIDERAÇÕES FINAIS (itens 6.5 e 18.6)
% Contribuições, impactos na tomada de decisão, coerência entre
% diagnóstico e proposta, valor gerado e perspectivas futuras.
% NÃO fazer: repetir problema ou metodologia, reapresentar dados e
% resultados, copiar trechos de seções anteriores.
% Esta seção NÃO deve conter citações de referências.
% =====================================================================
\section{Considerações Finais}

\pend{Contribuições práticas, valor gerado e perspectivas futuras. Sem
citações, sem reapresentar dados.}

% =====================================================================
% REFERÊNCIAS (itens 18.7, 19, 20 e 21)
% Alinhadas à esquerda, sem recuo, espaçamento simples. Apenas obras
% citadas no texto. Citação no estilo "autor, ano" ou "nome (ano)".
% Somente citações indiretas; apud não é permitido; nada de citação
% direta. Sigla entre COLCHETES na primeira ocorrência.
% Os modelos de cada tipo de fonte estão no item 21 do manual --- CONFIRA
% cada entrada contra o modelo correspondente antes de depositar.
% =====================================================================
\section{Referências}

\begin{referencias}
\refitem{FOIDL, H.; GOLENDUKHINA, V.; RAMLER, R.; FELDERER, M. Data
pipeline quality: influencing factors, root causes of data-related
issues, and processing problem areas for developers. \textbf{Journal of
Systems and Software}, v. 207, 111855, 2024.}

\refitem{\pend{demais referências --- conferir cada uma contra os modelos
do item 21 do manual, que distingue periódicos, relatórios de dados
online, livros, trabalhos acadêmicos, documentos jurídicos e anais.}}
\end{referencias}

% =====================================================================
% APÊNDICES E ANEXOS (opcionais, item 6.7)
% Apêndice: elaborado pelo autor. Anexo: documento externo.
% CONTAM para o limite de páginas.
% =====================================================================
% \newpage
% \section{Apêndice A --- Catálogo completo dos mecanismos de defeito}
% \pend{conteúdo}

\end{document}
